% Define styles
\tikzset{
    % Input/Output blocks
    io/.style={
        rectangle, 
        draw=black, 
        fill=blue!20,
        text width=3cm, 
        text centered, 
        minimum height=1cm,
        font=\footnotesize,
        drop shadow
    },
    % Convolution blocks
    conv/.style={
        rectangle, 
        draw=black, 
        fill=green!20,
        text width=2.8cm, 
        text centered, 
        minimum height=0.8cm,
        font=\footnotesize
    },
    % Residual block container
    resblock/.style={
        rectangle, 
        draw=black, 
        thick,
        dashed,
        fill=yellow!10,
        minimum width=3.5cm,
        minimum height=4.5cm,
        text width=3cm,
        align=center
    },
    % Small operation blocks
    operation/.style={
        rectangle, 
        draw=black, 
        fill=orange!20,
        text width=2.2cm, 
        text centered, 
        minimum height=0.6cm,
        font=\scriptsize
    },
    % Arrow style
    arrow/.style={
        ->,
        >=Stealth,
        thick
    },
    % Skip connection style
    skip/.style={
        ->,
        >=Stealth,
        thick,
        red,
        dashed
    }
}
\begin{tikzpicture}[htbp][node distance=1.5cm and 2cm]
    
    % Title
    \node[font=\Large\bfseries] at (7, 9) {Model Architecture Block Diagram};
    
    % Input
    \node[io] (input) at (0,0) {Input\\40-dim features\\20 timesteps};
    
    % Input Projection
    \node[conv, right=1.5cm of input] (proj_in) {1D Conv\\40 → 64 channels};
    
    % Draw 3 Temporal Blocks
    \foreach \i/\dil in {1/1, 2/2, 3/4} {
        \pgfmathsetmacro{\xpos}{3.5 + (\i-1)*4.5}
        
        % Residual block container
        \node[resblock] (resblock\i) at (\xpos, 0) {};
        
        % Label at top of residual block
        \node[font=\footnotesize\bfseries] at (\xpos, 1.8) {Temporal Block \i};
        \node[font=\tiny] at (\xpos, 1.5) {Dilation: \dil};
        
        % Components inside temporal block
        \node[operation] (conv1_\i) at (\xpos, 0.8) {Dilated Causal Conv\\kernel=2, dil=\dil};
        \node[operation] (wn1_\i) at (\xpos, 0.2) {Weight Norm};
        \node[operation] (relu1_\i) at (\xpos, -0.4) {ReLU};
        \node[operation] (drop1_\i) at (\xpos, -1.0) {Dropout (0.2)};
        \node[operation] (conv2_\i) at (\xpos, -1.6) {Dilated Causal Conv\\kernel=2, dil=\dil};
        
        % Internal connections
        \draw[arrow, line width=0.5pt] (conv1_\i) -- (wn1_\i);
        \draw[arrow, line width=0.5pt] (wn1_\i) -- (relu1_\i);
        \draw[arrow, line width=0.5pt] (relu1_\i) -- (drop1_\i);
        \draw[arrow, line width=0.5pt] (drop1_\i) -- (conv2_\i);
    }
    
    % Connect input to first temporal block
    \draw[arrow] (proj_in) -- ($(resblock1.west)+(0,0.8)$);
    
    % Connect temporal blocks sequentially
    \foreach \i in {1,...,2} {
        \pgfmathtruncatemacro{\next}{\i + 1}
        \draw[arrow] ($(resblock\i.east)+(0,-1.6)$) -- ($(resblock\next.west)+(0,0.8)$);
    }
    
    % Draw skip connections
    \foreach \i in {1,...,3} {
        \pgfmathsetmacro{\xpos}{3.5 + (\i-1)*4.5}
        % Skip connection path
        \draw[skip] ($(resblock\i.west)+(0,0.8)$) -- 
                     ($(resblock\i.west)-(0.3,0.8)$) -- 
                     ($(resblock\i.west)-(0.3,1.6)$) --
                     ($(resblock\i.east)+(0.3,1.6)$) --
                     ($(resblock\i.east)+(0.3,-1.6)$) --
                     ($(resblock\i.east)+(0,-1.6)$);
        
        % Add plus symbol for skip connection
        \node[circle, draw=red, thick, inner sep=1pt] at ($(resblock\i.east)+(0.15,-1.6)$) {+};
    }
    
    % Output Projection
    \node[conv, right=1.5cm of resblock3] (proj_out) {1×1 Conv\\64 → 5 channels};
    
    % Output
    \node[io, right=1.5cm of proj_out] (output) {Output\\5-dim sensor\\predictions};
    
    % Final connections
    \draw[arrow] ($(resblock3.east)+(0,-1.6)$) -- (proj_out);
    \draw[arrow] (proj_out) -- (output);
    
    % Legend
    \node[font=\footnotesize\bfseries] at (7, -3.5) {Legend:};
    \draw[skip] (5.5, -4) -- (6.5, -4) node[right, font=\tiny] {Skip Connection};
    \draw[arrow] (5.5, -4.3) -- (6.5, -4.3) node[right, font=\tiny] {Data Flow};
    
    % Architecture details box
    \node[draw, rectangle, fill=gray!10, text width=5cm, font=\tiny] at (7, -5.5) {
        \textbf{Architecture Details:}\\
        • 3 temporal blocks with skip connections\\
        • Exponential dilation: 1, 2, 4\\
        • Each block: 2 conv (k=2) + weight norm + ReLU\\
        • Dropout p=0.2 (regularization + MC inference)\\
        • 64 channels per block
    };
    
\end{tikzpicture}
